\chapter{Análise Crítica e Unificação}
\label{cha:critica}

\section{Crítica do trabalho atual}

A análise do projeto disponível no repositório revela que a principal limitação da versão anterior não estava na escolha do tema, mas na ausência de uma delimitação científica suficientemente operacional. O arquivo principal ainda apresentava grande quantidade de conteúdo genérico do template de TCC, enquanto a discussão específica sobre metaverso, critérios de seleção de literatura, comparação de estudos e perguntas de pesquisa ainda não estava consolidada.

Outro ponto crítico é a tendência de reunir em uma única definição elementos de naturezas diferentes. Persistência, interação por avatares e existência de um ambiente tridimensional descrevem propriedades do fenômeno. Já realidade virtual, inteligência artificial, blockchain e computação em nuvem descrevem tecnologias que podem viabilizar essas propriedades. Misturar essas camadas pode tornar a definição circular ou excessivamente dependente de tecnologias específicas.

A terceira limitação é metodológica. Uma tabela PRISMA só é cientificamente válida como registro quantitativo quando os números correspondem a uma busca realmente executada e documentada. Por isso, nesta versão os números apresentados são explicitamente identificados como pertencentes à busca exploratória comparativa, e não como uma revisão sistemática completa.

\section{Comparação conceitual}

A definição do projeto e o estudo de Tukur et al. (2024) apresentam pontos de convergência, mas partem de objetivos diferentes. O projeto procura caracterizar o que o usuário vivencia e quais propriedades o ambiente deve apresentar. Tukur et al. procuram identificar as técnicas e tecnologias que sustentam a construção desses ambientes.

Ritterbusch e Teichmann (2023) fornecem uma terceira camada de validação conceitual ao mostrar que características como ambientes tridimensionais, avatares, interação, imersão e economia virtual aparecem de forma recorrente na literatura, embora não exista consenso absoluto sobre uma única definição.

Dessa comparação, resulta uma conclusão metodológica: a unificação é adequada desde que seja feita por camadas. Não se recomenda simplesmente juntar todas as características em uma lista única e tratá-las como requisitos obrigatórios.

\section{Definição unificada}

Com base na comparação realizada, este trabalho adota a seguinte definição operacional:

\begin{quote}
\textit{O metaverso é um ecossistema de ambientes digitais tridimensionais e compartilhados, nos quais usuários representados por identidades digitais podem interagir entre si e com objetos virtuais em experiências imersivas ou espacialmente contextualizadas. Esses ambientes podem apresentar persistência, interoperabilidade entre espaços, mecanismos de criação e circulação de ativos digitais e formas de atividade econômica, sendo sustentados por um conjunto variável de tecnologias, como realidade estendida, inteligência artificial, sistemas de comunicação, computação distribuída e tecnologias descentralizadas.}
\end{quote}

Essa definição preserva os elementos centrais da proposta inicial e incorpora a perspectiva tecnológica identificada na literatura. A palavra ``podem'' é deliberada: persistência, interoperabilidade e economia digital são características relevantes para a definição operacional deste trabalho, mas não devem ser tratadas como requisitos tecnológicos obrigatórios de toda implementação existente.

\section{Concordância com a unificação}

A unificação é considerada adequada para o objetivo deste projeto porque resolve uma diferença de escala entre as abordagens comparadas. A definição inicial descreve principalmente a experiência e as propriedades desejadas; Tukur et al. descrevem os meios tecnológicos e os contextos de aplicação; Ritterbusch e Teichmann ajudam a delimitar as características conceituais recorrentes.

A concordância, contudo, é parcial no sentido metodológico: não se deve afirmar que a definição unificada representa a definição científica definitiva do metaverso. Ela representa uma definição operacional deste trabalho, construída a partir da literatura analisada.

Essa distinção fortalece o texto acadêmico porque permite apresentar uma posição clara sem transformar uma síntese específica em consenso inexistente na literatura.